\documentclass[11pt]{article}
\usepackage[margin=0.9in]{geometry}
\usepackage{booktabs}
\usepackage{array}
\usepackage{tabularx}
\usepackage{longtable}
\usepackage{hyperref}
\hypersetup{colorlinks=true, linkcolor=blue, urlcolor=blue}
\setlength{\parindent}{0pt}
\setlength{\parskip}{0.55em}
\begin{document}

{\Large \textbf{Prompt-Profile RFM Third-Stage Steering Grand Plan}}\\
Generated: 2026-04-23 23:28 UTC\\
Canonical repo note:\quad \texttt{docs/prompt-profile-rfm-third-stage-steering-plan.md}

\section*{Object}
The third-stage steering object is mode-local. A scientific row must close this
full chain inside one \texttt{(dataset, thinking mode)} path:
\begin{center}
\begin{tabular}{c}
\texttt{stats collection -> prompt-level materialization}\\
\texttt{-> probe/RFM training -> vector export -> steering}
\end{tabular}
\end{center}
Thinking-on and thinking-off
are separate experimental objects. Cross-mode steering rows and mixed-mode
average vectors are not evidence for this stage.

\section*{Active Data Foundation}
The active data foundation is the rebuilt four-dataset rollout-stat suite under
\texttt{Qwen/Qwen3-1.7B}, with \texttt{temperature=0.2},
\texttt{num\_generations=10}, and the max-token / max-model-length contract
recorded in the collector manifests. Both thinking modes are collected for every
dataset.

\begin{center}
\begin{tabular}{lrl}
\toprule
Dataset & Active size & Thinking modes \\
\midrule
\texttt{LiveCodeBench} & full \texttt{1055} & \texttt{on}, \texttt{off} \\
\texttt{TACO-hard} & \texttt{1000 / 5536} & \texttt{on}, \texttt{off} \\
\texttt{MATH level-5} & \texttt{1000 / 2304} & \texttt{on}, \texttt{off} \\
\texttt{Omni-MATH >= 7} & full \texttt{916} & \texttt{on}, \texttt{off} \\
\bottomrule
\end{tabular}
\end{center}

\texttt{LiveCodeBench-extra} is not a separate steering benchmark; it is a
subset of \texttt{LiveCodeBench release\_v6}. The current stage starts from the
new rebuilt mode-local archives, not from any earlier steering bundle.

\section*{Pipeline}
\textbf{P0: stats receipt.} Each archive must carry dataset/filter, thinking
mode, prompt formatter, model/tokenizer revision where available, generation
config, prompt ids, prompt text, prompt token ids, completion text and token ids,
correctness where available, length metrics, loop flags, cap-hit flags, and raw
row metadata.

\textbf{P1: prompt-level materialization.} Recompute labels from the same
mode-local archive. Required targets are \texttt{majority\_s\_0.5},
\texttt{mean\_relative\_length}, \texttt{p\_loop}, and \texttt{p\_cap} when
available. \texttt{majority\_s\_0.5} means strict-majority long-rollout risk
under \texttt{relative\_length >= 0.5}; it is not a pure loop label.

\textbf{P2: detector table.} For each admitted object, compare prompt-only
baselines, activation-linear, activation-MLP, and layerwise RFM on the same
mode-local splits. Primary detector metric is \texttt{PR-AUC}; \texttt{ROC-AUC}
and threshold metrics are diagnostic.

\textbf{P3: vector export.} Export one signed vector per transformer block from
the mode-local RFM bundle. \texttt{+v\_l} is the higher-risk direction at block
\texttt{l}; \texttt{-v\_l} is the anti-risk direction. Before interpreting
steering, report vector hashes, sign rule, projection scores, bootstrap/seed
cosine stability, and cross-layer cosine structure.

\textbf{P4: steering.} Run steering only after the corresponding detector and
vector receipts exist for the same \texttt{(dataset, thinking mode)} object.
\texttt{LiveCodeBench} is first because the grader and code-generation
evaluation path are the cleanest.

\section*{Admission And Direction Gates}
A dataset/mode is admitted into steering only if its mode-local train or
fit-train object has at least \texttt{10\%} positives under
\texttt{majority\_s\_0.5}. Sparse-positive objects remain diagnostic unless
explicitly accepted.

For direction quality, fewer than \texttt{20 / 28} layers clearing mean
bootstrap cosine \texttt{>= 0.7} means the steering run may still be useful as
infrastructure, but not as a stable-vector causal read without a caveat. If the
selected detector layer is unstable, report it directly.

\section*{Fixed Steering Contract}
\begin{center}
\begin{tabularx}{\textwidth}{lX}
\toprule
Field & Contract \\
\midrule
Timing & Prefill forward pass only \\
Hook site & \texttt{prefill\_layer\_output\_all\_tokens} \\
Token scope & Every prompt token, not only the last prompt token \\
Block scope & Every block uses its own block-specific vector \texttt{v\_l} \\
Linear strength & \texttt{epsilon = 0.2} \\
Spherical strength & \texttt{t = 0.3} \\
Decode budget & Full source-manifest budget; reduced caps are smoke-only \\
\bottomrule
\end{tabularx}
\end{center}

First-pass conditions:
\begin{center}
\begin{tabular}{ll}
\toprule
Condition & Meaning \\
\midrule
\texttt{no\_steer} & matched baseline \\
\texttt{minus\_v\_linear} & additive anti-risk direction \\
\texttt{plus\_v\_linear} & additive sign-flip / risk direction \\
\texttt{random\_linear} & additive random-direction control \\
\texttt{minus\_v\_spherical} & norm-preserving anti-risk pole \\
\texttt{plus\_v\_spherical} & norm-preserving sign-flip pole \\
\texttt{random\_spherical} & norm-preserving random-direction control \\
\bottomrule
\end{tabular}
\end{center}

Do not add layer selection, probe gating, online controllers, or strength sweeps
until the seven-condition table exists for the relevant mode-local object.

\section*{Required Outputs}
Each steering run must write a condition summary, prompt-level ledger, and a
steering-run v1 record with prompt ids,
prompt-id hash, vector artifact hash, thinking mode, hook site, \texttt{t},
\texttt{epsilon}, generation config, grader version, git commit, and output
root.

Headline metrics are \texttt{pass@1}, accuracy delta vs \texttt{no\_steer},
loop fraction, max-length-hit fraction, long-rollout fraction, average length,
median length, and bootstrap intervals for key deltas when sample size supports
them. Intervention diagnostics must include pre/post norm, spherical
norm-preservation error, starting angle, realized angular movement, condition
runtime, and timeout or cap-hit asymmetry.

\section*{Execution Backlog}
\begin{enumerate}
\item Finish all eight rebuilt stats archives for the four datasets and two
thinking modes.
\item Materialize mode-local prompt-profile objects and admission-gate tables.
\item Train the \texttt{LiveCodeBench} thinking-on and thinking-off detector
surfaces from the rebuilt archives.
\item Export the two \texttt{LiveCodeBench} mode-local vector bundles and
direction diagnostics.
\item Run the full seven-condition \texttt{LiveCodeBench} table in thinking
\texttt{on}.
\item Run the full seven-condition \texttt{LiveCodeBench} table in thinking
\texttt{off}.
\item Promote \texttt{TACO-hard}, \texttt{MATH level-5}, and
\texttt{Omni-MATH >= 7} path-by-path only after their rebuilt mode-local
materializations pass the gate and have detector/vector receipts.
\item Open transfer only after at least two non-identical datasets inside the
same thinking mode have admitted detector/vector bundles, direction diagnostics,
and readable in-distribution steering receipts.
\end{enumerate}

\section*{Positive Claim Gate}
A steering row can be called positive only when the matched \texttt{minus\_v}
condition improves loop, cap-hit, or long-rollout metrics relative to
\texttt{no\_steer}; beats the matched random control; is not matched by the
sign-flip condition; does not incur an unacceptable accuracy drop; uses full
held-out prompts and the full decode contract; and matches a mode-local,
benchmark-local vector bundle by prompt ids and vector hashes.

If those conditions fail, report the result as a negative steering read,
implementation receipt, or diagnostic control, not as evidence that steering
works.

\section*{Exclusions}
Do not use vector bundles outside the rebuilt mode-local path. Do not claim
\texttt{LiveCodeBench-extra} as a separate benchmark. Do not call
\texttt{majority\_s\_0.5} a pure loop label. Do not treat detector
\texttt{PR-AUC} as a steering result. Do not treat reduced-cap, prompt-subset,
or last-token-only rows as full-stage evidence. Do not mix thinking-on and
thinking-off vectors in transfer.

\end{document}
